\Needspace{7\baselineskip}
\section{Relational Geometry: The FCC/RD Cell}
\label{sec:04}
Figure \(\ref{fig:fcc-rd-cell}\) gives a coordinate presentation of the derived cell geometry and its centered maximal inball.  The diagram is a bookkeeping aid: the relational derivation precedes the drawing, and neither the drawing nor a background lattice supplies scientific authority.

Section 3 closed the base book. It fixed the unit, the closure traversal, the closure-service support, the connectivity-channel count, the channel distribution, and the closure-bookkeeping value. Section 4 asks how that closed book can be represented geometrically without importing a background space as an additional primitive.

The answer is a specific relational geometry: the face-centered-cubic lattice and its rhombic-dodecahedral cell. The claim is not that physical space has been assumed to be a lattice. The claim is that, once the relational bookkeeping has three connectivity channels and a two-channel transfer edge, the allowed two-active, one-idle orbit determines a twelve-edge relational shell. The integer lattice generated by that shell is the FCC presentation, and its Voronoi cell is the rhombic dodecahedron.

\Needspace{5\baselineskip}
\subsection{Presentation, Not Background}
The coordinate triples used in this section are a presentation layer. They let the paper compute exactly, but they are not an observer-independent stage on which existence is placed. The three coordinate slots are tied to the three connectivity channels by the seam theorem: the channel-coupling action and the FCC-shell action are the same transitive relational action. The coordinate picture is therefore a drawing of the relational structure, not an added physical grid.

This matters because the same numeral can appear for different reasons. The symbol \(c\) denotes the connectivity-channel count, \(c=3\). It is not silently identified with spatial dimension. When a three-coordinate model appears here, it appears under the geometry presentation license, and every norm, radius, volume, and cell construction remains internal to the presentation until a later bridge gives it physical-unit interpretation.

A coordinate component written \(0\) in this section is not a zero quantity and is not a held value. It marks channel absence in the presentation tuple: the signed-entry rule is applied only to nonzero components, so no \((\sigma,0)\) pair and no zero magnitude is admitted. Thus tuples such as \((\pm1,\pm1,0)\), \((0,1,1)\), and labels such as \(\operatorname{Vor}(0)\) use \(0\) only as presentation notation or a reference label, never as an object in the positive quantity domain.

% LAYOUT_ONLY_GUARD:GEOMETRY_FORCING_CHAIN
\Needspace{17\baselineskip}
In this section, "forced" has a precise local meaning: determined up to channel relabeling, sign convention, and uniform normalization by the allowed two-active, one-idle transfer orbit. The chain is:

\begin{equation*}
\begin{adjustbox}{max width=\linewidth,center}
$\displaystyle
\text{three channels}
\rightarrow
\text{two-active/one-idle transfer edge}
\rightarrow
\text{closure under relabeling/sign}
\rightarrow
\text{unique twelve-vector orbit}
\rightarrow
\text{lattice generated by that orbit}
\rightarrow
\text{even-sum integer lattice}
\rightarrow
\text{FCC presentation}
\rightarrow
\text{RD Voronoi cell}.
$
\end{adjustbox}
\end{equation*}
The FCC lattice is therefore not chosen as a background lattice. It is the forced integer lattice generated by the required twelve-edge orbit.

% LAYOUT_ONLY_GUARD:EDGE_COORDINATE_SENTENCE
\Needspace{7\baselineskip}
The seam result also fixes what an edge is in this geometry. A transfer requires two distinct registers: a source and a target. With three channels available, an edge couples two channels while the third is idle. The unique size-twelve orbit of the model relabeling group is exactly the two-active, one-idle pattern. In coordinates this is the \((\pm 1,\pm 1,0)\)-type orbit, but that coordinate form is the presentation of the relational fact: an edge is the minimal two-channel coupling.

The three-active \((a,a,a)\)-type orbit is therefore not an edge candidate: it is the closure/cycle object, so the eight-direction BCC shell is excluded on transfer-typing grounds, not by a norm choice.

\Needspace{5\baselineskip}
\subsection{The FCC Lattice and Its Minimal Shell}
The integer lattice generated by that required twelve-edge orbit is

\begin{equation*}
\begin{adjustbox}{max width=\linewidth,center}
$\displaystyle
\Lambda_{\mathrm{FCC}}
:=
\operatorname{span}_{\mathbb Z}\{(0,1,1),(1,0,1),(1,1,0)\}
=
\{(x,y,z)\in\mathbb Z^3:\ x+y+z\ \text{is even}\}.
$
\end{adjustbox}
\end{equation*}
The two descriptions agree exactly. Each generator has even coordinate sum, and every even-sum integer triple can be written in the generator basis by

\begin{equation*}
\begin{adjustbox}{max width=\linewidth,center}
$\displaystyle
n_1=\frac{-x+y+z}{2},\qquad
n_2=\frac{x-y+z}{2},\qquad
n_3=\frac{x+y-z}{2}.
$
\end{adjustbox}
\end{equation*}
% LAYOUT_ONLY_GUARD:STANDALONE_THUS
\Needspace{7\baselineskip}
Thus

\begin{equation*}
\begin{adjustbox}{max width=\linewidth,center}
$\displaystyle
(x,y,z)=n_1(0,1,1)+n_2(1,0,1)+n_3(1,1,0).
$
\end{adjustbox}
\end{equation*}
The generator determinant has absolute value \(2\), so this is a rank-three sublattice of index \(2\) and covolume \(2\). This display identifies the generated presentation object exactly; it does not turn the lattice into a background space.

% LAYOUT_ONLY_GUARD:FCC_MINIMAL_SHELL
\Needspace{12\baselineskip}
Descriptively, once the two-active orbit has already forced the shell, that same shell is the minimal shell of the FCC lattice. The native geometry top-invariant theorem identifies \(\eta\) as the geometry column's top invariant. Every nonzero vector in \(\Lambda_{\mathrm{FCC}}\) has squared norm at least \(2\), and equality occurs exactly for the twelve vectors with two coordinates \(\pm 1\) and one coordinate \(0\):

\begin{equation*}
\begin{adjustbox}{max width=\linewidth,center}
$\displaystyle
S_{12}
=
\{v\in\Lambda_{\mathrm{FCC}}:\|v\|^2=2\}
=
\{\text{permutations and sign patterns of }(1,1,0)\},
\qquad |S_{12}|=12 .
$
\end{adjustbox}
\end{equation*}
The nearest-neighbor edge invariant is

\begin{equation*}
\begin{adjustbox}{max width=\linewidth,center}
$\displaystyle
\eta=\sqrt{2}.
$
\end{adjustbox}
\end{equation*}
Thus \(\eta\) plays in the geometry column the role that \(I\) and \(c\) play in the first two columns: by the Section 3 invariant theorem, \(I\) and \(c\) are the booked-to top entries of \(M_{\mathrm{base}}\); by the geometry top-invariant theorem, \(\eta\) is the geometry column's top invariant. It singles out no channel and no absolute scale, with the lower entries carrying it through the cell-volume trace.

This value is not assigned. It is the minimal nonzero norm of the required twelve-edge orbit in the identified FCC presentation. The \(\sqrt{2}\) that later appears in the cell volume, filled share, and residual share is this same edge invariant carried through the geometry, not a new constant and not a power law.

This is the same primitive-versus-representation distinction already met by the closure quantities. The exact forms \(2\pi\), \(1/(2\pi)\), and \(\sqrt2\) belong to established closure, support, and norm objects. Their appearance in the displayed arrays does not enlarge the primitive ontology, and it does not license arbitrary exact expressions without their own derivation and typing.

\Needspace{5\baselineskip}
\subsection{The Rhombic-Dodecahedral Cell}
The cell is built from the twelve minimal shell directions. For a normalization parameter \(u>0\), define

\begin{equation*}
\begin{adjustbox}{max width=\linewidth,center}
$\displaystyle
C_u
=
\bigcap_{n\in S_{12}}\{p\in\mathbb R^3:\ n\cdot p\le 2u\}
=
\{(x,y,z): |x\pm y|\le 2u,\ |x\pm z|\le 2u,\ |y\pm z|\le 2u\}.
$
\end{adjustbox}
\end{equation*}
Volume enters here only after the relational shell has forced a coordinate presentation. All volumes in this section are presentation volumes induced by the licensed coordinate model. They compare derived objects inside the same relational presentation: the FCC/RD cell, the lattice covolume, and the inradius ball. They are not physical volumes and carry no physical-unit interpretation.

Solving this finite half-space system gives the complete cell census. The cell has exactly twelve congruent rhombic facets, one for each shell direction; fourteen vertices, split as

\begin{equation*}
\begin{adjustbox}{max width=\linewidth,center}
$\displaystyle
V_3=\{(\pm u,\pm u,\pm u)\},\qquad
V_4=\{(\pm 2u,0,0),(0,\pm 2u,0),(0,0,\pm 2u)\};
$
\end{adjustbox}
\end{equation*}
six antipodal facet pairs; and exact model-internal volume

\begin{equation*}
\begin{adjustbox}{max width=\linewidth,center}
$\displaystyle
\operatorname{Vol}(C_u)=16u^3.
$
\end{adjustbox}
\end{equation*}
This polytope is the rhombic dodecahedron. The name is conventional; the proof does not rely on the name. The shape is obtained by enumerating the constraints and vertices directly.

The same cell is also the Voronoi cell of \(\Lambda_{\mathrm{FCC}}\). Farther lattice vectors do not cut the cell: vectors of squared norm \(4\) only touch at the axis vertices, and all larger shells are slack. Consequently

\begin{equation*}
\begin{adjustbox}{max width=\linewidth,center}
$\displaystyle
\operatorname{Vor}_{\Lambda_{\mathrm{FCC}}}(0)=C_{1/2},
$
\end{adjustbox}
\end{equation*}
% LAYOUT_ONLY_GUARD:CELL_VOLUME
\Needspace{9\baselineskip}
and the translates \(\xi+C_{1/2}\), \(\xi\in\Lambda_{\mathrm{FCC}}\), tile the presentation space with disjoint interiors. For \(C_{1/2}\), the Voronoi cell volume is the lattice covolume, \(2\). In inradius form, with \(r\) the model-internal distance from the cell center to a facet,

\begin{equation*}
\begin{adjustbox}{max width=\linewidth,center}
$\displaystyle
\operatorname{Vol}(\operatorname{Cell}(r))=4\sqrt{2}\,r^3 .
$
\end{adjustbox}
\end{equation*}
The proof uses two agreeing routes: the lattice covolume route and the direct \(16u^3\) route. Both are exact and symbolic.

\Needspace{5\baselineskip}
\subsection{Filled Share and Residual Share}
Place the inradius ball \(B(r)\) inside each cell. It touches all twelve facet planes, and neighboring inradius balls are tangent along the face-neighbor directions. The filled share of the cell is then the ball volume divided by the cell volume:

\begin{equation*}
\begin{adjustbox}{max width=\linewidth,center}
$\displaystyle
D
=
\frac{\operatorname{Vol}(B(r))}{\operatorname{Vol}(\operatorname{Cell}(r))}
=
\frac{\tfrac{4}{3}\pi r^3}{4\sqrt{2}\,r^3}
=
\frac{\pi}{3\sqrt{2}}.
$
\end{adjustbox}
\end{equation*}
The normalization cancels. No global packing theorem is used or needed. The value is the filled share of this cell configuration, not a claim about optimal sphere arrangements and not a physical matter-density claim.

\Needspace{5\baselineskip}
\subsection{Unique Centered Inball and Dual Geometric Share}
The centered inradius ball is not selected from arbitrary subregions. For a cell written as the intersection of antipodal supporting half-spaces

\begin{equation*}
\begin{adjustbox}{max width=\linewidth,center}
$\displaystyle
n\cdot p\le r,
\qquad
-n\cdot p\le r,
$
\end{adjustbox}
\end{equation*}
let a candidate inscribed ball have center \(z\) and radius \(R\). The two inequalities imply

\begin{equation*}
\begin{adjustbox}{max width=\linewidth,center}
$\displaystyle
n\cdot z+R\le r,
\qquad
-n\cdot z+R\le r.
$
\end{adjustbox}
\end{equation*}
Adding them gives \(R\le r\). Equality forces \(n\cdot z=0\) for every facet normal, and those normals span the presentation space, so \(z=0\). Thus the centered ball \(B(r)\) is the unique maximal inscribed ball.

Its boundary area and the cell boundary area are

\begin{equation*}
\begin{adjustbox}{max width=\linewidth,center}
$\displaystyle
\operatorname{Area}(\partial B(r))=4\pi r^2,
\qquad
\operatorname{Area}(\partial\operatorname{Cell}(r))
=12\sqrt2\,r^2.
$
\end{adjustbox}
\end{equation*}
Consequently the volume and boundary-area ratios agree exactly:

\begin{equation*}
\begin{adjustbox}{max width=\linewidth,center}
$\displaystyle
\boxed{
\frac{\operatorname{Vol}(B(r))}
{\operatorname{Vol}(\operatorname{Cell}(r))}
=
\frac{\operatorname{Area}(\partial B(r))}
{\operatorname{Area}(\partial\operatorname{Cell}(r))}
=
\frac{\pi}{3\sqrt2}
=D.
}
$
\end{adjustbox}
\end{equation*}
The equality follows from the common tangential-polytope identity \(V=rA/3\). It makes \(D\) a scale-free geometric invariant of the centered-inball/full-cell relation. It does not make the inball boundary a subset of the cell boundary, and the complementary scalar \(\Omega=I-D\) is not a boundary patch or substance.

The complementary cell share is recorded in transfer form:

\begin{equation*}
\begin{adjustbox}{max width=\linewidth,center}
$\displaystyle
D+\Omega=I,
\qquad
\Omega=1-\frac{\pi}{3\sqrt{2}},
\qquad
0<\Omega<1.
$
\end{adjustbox}
\end{equation*}
The residual is a positive booked share of the same cell unit. It is not dark matter, not mass, and not a measured physical density. If a later section assigns additional physical meaning to \(\Omega\), it must do so through its own bridge or identity claim. Section 4 only supplies the exact geometry and the exact residual bookkeeping.

\Needspace{5\baselineskip}
\subsection{The Cell Matrix}
For a compact presentation, Section 4 appends one column to the base display of Section 3. The grouped geometry entries are

\begin{equation*}
\begin{adjustbox}{max width=\linewidth,center}
$\displaystyle
\begin{pmatrix}
\eta\\[2pt]
D\\[2pt]
\Omega
\end{pmatrix}
=
\begin{pmatrix}
\sqrt{2}\\[2pt]
\dfrac{\pi}{3\sqrt{2}}\\[2pt]
1-\dfrac{\pi}{3\sqrt{2}}
\end{pmatrix}.
$
\end{adjustbox}
\end{equation*}
The top entry \(\eta\) is the edge invariant, just as \(c\) is the fair connectivity-channel invariant in the second column. The entries \(D\) and \(\Omega\) carry that same edge through the cell-volume trace. This is an identification by provenance, not a functional law \(D=f(\eta)\). The geometry has one forced edge invariant and two exact shares carrying it.

Thus the cell matrix is

\begin{equation*}
\begin{adjustbox}{max width=\linewidth,center}
$\displaystyle
M_{\mathrm{cell}}
=
\begin{pmatrix}
I & c & \eta\\[2pt]
\tau & E & D\\[2pt]
m & G & \Omega
\end{pmatrix}
=
\begin{pmatrix}
1 & 3 & \sqrt{2}\\[2pt]
2\pi & \dfrac{9}{2\pi} & \dfrac{\pi}{3\sqrt{2}}\\[2pt]
\dfrac{1}{2\pi} & 18\pi & 1-\dfrac{\pi}{3\sqrt{2}}
\end{pmatrix}.
$
\end{adjustbox}
\end{equation*}
The display appends this third column while preserving every base entry. The arrangement records their independently established provenance; it does not derive the values or license matrix algebra.

\Needspace{5\baselineskip}
\subsection{Static Cell Referent}
The tiling also gives a static referent for later use. Each lattice site \(\xi\) determines exactly one cell, and the cell determines its site:

\begin{equation*}
\begin{adjustbox}{max width=\linewidth,center}
$\displaystyle
\operatorname{REF}(\xi):=(\xi,\xi+C_{1/2}).
$
\end{adjustbox}
\end{equation*}
This is only a site-cell referent supplied by the geometry. Runtime questions about posting, ordering, and rates belong to the later growth construction.

Section 4 therefore completes the displayed geometry group without crossing the bridge to physical units. It gives the FCC/RD relational realization, the exact cell, the exact filled and residual shares, and the cell matrix.

It has licensed only the static site-cell referent and the geometry results collected in \(M_{\mathrm{cell}}\). Thus \(M_{\mathrm{cell}}\) closes only the static relational geometry: the FCC/RD presentation, the cell, the edge invariant, and the filled/residual shares. It does not install a background grid, does not make cells into observer systems, and does not introduce observer orientation, observer distinguishability, posting order, or runtime growth. At this stage no cell has been typed as an observer system, no observer-orientation or observer-distinguishability structure has been introduced, and no posting dynamics have begun. Those are later constructions, not hidden assumptions inside \(M_{\mathrm{cell}}\). It does not promote FCC/RD to ontology, does not treat model-internal distances as measured lengths, does not derive SI units, does not read \(\Omega\) as mass, and does not turn geometry counts into earlier bookkeeping reasons.

This layer can also lose, and says how. Exhibit an admissible two-active, one-idle transfer orbit that does not force the twelve-member \((a,a,0)\) shell, and the forced FCC/RD route falls. Exhibit a vector in that shell with odd coordinate sum, or an even-sum FCC generator not reached by the shell, and the lattice identification falls. Exhibit a lattice vector outside the minimal shell that cuts \(\operatorname{Vor}_{\Lambda_{\mathrm{FCC}}}(0)\), and the rhombic-dodecahedral cell claim falls. Exhibit a dependence of \(D\) or \(\Omega\) on the normalization \(r\), and the filled/residual share bookkeeping falls. Exhibit a coordinate \(0\) that enters as a zero quantity rather than channel absence in the presentation tuple, and the presentation discipline falls. Each failure mode is specific and checkable.

